\chapter{Literature Review}

\section{Theoretical Foundation}
\label{sec:theory}

\subsection{Agency Theory}

According to \textcite{jensen1976}, owners appoint managers to run the firm on their behalf, but managers have more information about the firm than owners do, and the interests of the two parties may differ. Owners therefore incur monitoring costs. The external audit and the audit committee are two such monitoring mechanisms, as both exist to check the financial information that managers report. Where these mechanisms work effectively, managers should have less opportunity to manage earnings.

\subsection{Positive Accounting Theory}

\textcite{watts1986} argue that managers select accounting methods and estimates that serve their own contracting interests and those of the firm, for example, in relation to bonus plans and debt covenants. Since accruals rely heavily on judgment, they provide an obvious means of influencing reported results.

\subsection{Audit Quality}

\textcite{deangelo1981} defines audit quality as the probability that an auditor will both discover and report a breach in the client's accounting system. This definition has two dimensions: the auditor's competence, which determines whether the breach is found, and the auditor's independence, which determines whether it is reported. \textcite{becker1998} associate higher audit quality, proxied by auditor size, with lower discretionary accruals. However, this proxy performs poorly in Pakistan, where \textcite{abid2018} find no significant difference in earnings management between clients of Big Four and other auditors. For this reason, this study focuses on three characteristics that act on either competence or independence: audit tenure, audit fees and audit committee independence.

\subsection{Learning and Familiarity in the Auditor--Client Relationship}
\label{sec:learning}

Two competing views explain how tenure may affect audit quality. The learning view holds that a longer engagement allows the auditor to build client-specific knowledge, which improves the audit. The familiarity view, by contrast, holds that a long relationship may reduce professional skepticism and weaken independence. The two views therefore predict opposite effects of tenure. Under Regulation 33 of the 2019 Regulations, a non-financial listed company is only required to rotate the engagement partner every five years, not the audit firm itself. Audit firm tenure can therefore become lengthy, which makes it an important feature to examine in this sector.

\subsection{Audit Effort and Economic Bonding}
\label{sec:fees-theory}

A higher audit fee may pay for greater audit effort or may reflect a riskier client, and \textcite{simunic1980} shows that audit fees are determined by both. A fee that is unusually large in relation to the audit firm's overall income can also create economic dependence, which may make the auditor less willing to challenge management \parencite{choi2010}. The overall effect of audit fees on earnings management therefore cannot be predicted in advance.

\subsection{The Audit Committee as a Monitoring Mechanism}
\label{sec:committee-theory}

Independent directors have fewer reasons to side with management, so an independent audit committee should oversee financial reporting and the external audit more effectively. \textcite{klein2002} provides the main empirical support for this view.

\section{Measurement of Earnings Management}

Accrual-based earnings management uses the discretion that accrual accounting allows, while real earnings management changes actual business activities such as production or discretionary spending. This study measures the accrual-based form, following the review by \textcite{healy1999}.

\begin{description}[leftmargin=0pt,labelindent=0pt,style=unboxed]
  \item[Total accruals (TA):] net income minus operating cash flow, scaled by opening total assets.
  \item[Non-discretionary accruals (NDA):] the part of total accruals explained by the change in revenue, the change in receivables and property, plant and equipment.
  \item[Discretionary accruals (DA):] total accruals minus non-discretionary accruals. This is the part that managers control.
  \item[Earnings management (EM):] the absolute value of discretionary accruals, which captures the size of the manipulation whatever its direction.
\end{description}

The Modified Jones Model of \textcite{dechow1995} extends the \textcite{jones1991} model by removing the change in receivables from the revenue change when predicting non-discretionary accruals, so that credit sales are not treated as non-discretionary. \textcite{kothari2005} show that accrual measures can be affected by firm performance, which is why return on assets enters the model as a control variable.

\section{Audit Tenure and Earnings Management (H1)}

\textcite{johnson2002} find lower quality of financial reports in the early years of an engagement and no evidence of a decline in quality at longer tenure. \textcite{myers2003} report that longer tenure is associated with better earnings quality and conclude that their results do not support mandatory rotation. Pakistani evidence points the same way. \textcite{kalabeke2019}, using 280 non-financial PSX companies over 2008--2017, find that longer audit firm tenure is negatively associated with accrual-based earnings management. \textcite{alam2023} also report a negative association for cement companies.

The familiarity argument predicts the opposite, and it underlies proposals for mandatory rotation. Because the international evidence is divided and the Pakistani samples cover other periods and sectors, this study tests a two-sided hypothesis.

\emph{Theoretical basis:} the learning and familiarity views in Section~\ref{sec:learning}, both within agency theory, which describe how the auditor's monitoring of managers changes as the relationship lengthens.

\begin{quote}
  \textbf{H1:} Audit tenure is significantly associated with earnings management.
\end{quote}

\section{Audit Fees and Earnings Management (H2)}

If fees reflect audit effort, higher fees should go with less earnings management. \textcite{frankel2002} find that audit fees are negatively associated with several indicators of earnings management, while fees for non-audit services are positively associated with them. If higher fees instead reflect client risk or create economic dependence, higher fees should go with more earnings management. \textcite{choi2010} report that abnormally high audit fees go with lower audit quality, measured by absolute discretionary accruals. In Pakistan, \textcite{abid2018} point to strong economic bonding between auditors and their clients. Direct evidence from Pakistan on audit fees and accrual-based earnings management is limited.

Because the fee alone cannot separate effort from risk and dependence, no sign is predicted.

\emph{Theoretical basis:} the effort and economic bonding arguments in Section~\ref{sec:fees-theory}, which describe the audit fee as either strengthening or weakening the auditor's monitoring role.

\begin{quote}
  \textbf{H2:} Audit fees are significantly associated with earnings management.
\end{quote}

\section{Audit Committee Independence and Earnings Management (H3)}

\textcite{klein2002} finds a negative relation between audit committee independence and abnormal accruals. \textcite{bedard2004} report that audit committees made up of independent members are associated with less aggressive earnings management, and \textcite{xie2003} examine the role of the board and the audit committee in constraining earnings management.

Pakistani evidence is mixed but leans towards the monitoring view. \textcite{latif2015} find that audit committee independence is negatively associated with earnings management in non-financial firms listed on the Karachi Stock Exchange, and \textcite{khan2020} find that audit committee effectiveness reduces discretionary accruals in PSX-listed firms. \textcite{khan2024} find that board independence and financial experts on the board reduce discretionary accruals after the 2017 Code of Corporate Governance. In contrast, \textcite{mumtaz2018} reports that more independent audit committee members are linked with more earnings manipulation. A negative relationship is expected.

\emph{Theoretical basis:} agency theory and the monitoring role of the audit committee in Section~\ref{sec:committee-theory}, under which independent members reduce the scope for managers to manage earnings.

\begin{quote}
  \textbf{H3:} Greater audit committee independence is associated with lower earnings management.
\end{quote}

\section{Control Variables}
\label{sec:controls}

\begin{description}[leftmargin=0pt,labelindent=0pt,style=unboxed]
  \item[Firm size:] larger firms face more scrutiny but also have more complex accruals, so no sign is predicted.
  \item[Leverage:] indebted firms may manage earnings to avoid breaching loan covenants \parencite{defond1994}, and \textcite{kamran2014} find that Pakistani firms with higher leverage manage earnings more.
  \item[Return on assets:] accruals are related to performance \parencite{kothari2005}, and \textcite{kamran2014} find that more profitable Pakistani firms manage earnings more actively.
\end{description}

\section{Conceptual Framework and Model Development}

The framework tests the association of three audit-related variables with earnings management after controlling for company characteristics. Audit tenure, audit fees and audit committee independence are the independent variables because each is a measurable feature of the monitoring arrangements described in Section~\ref{sec:theory}. Earnings management is the outcome these arrangements are expected to constrain. Firm size, leverage and return on assets are controlled because earlier studies link them to accrual choices (Section~\ref{sec:controls}). Model~1 in Section~\ref{sec:model-spec} turns the framework into a regression equation.

\begin{table}[ht]
  \centering
  \caption{Conceptual framework}
  \label{tab:framework}
  \small
  \begin{tabular}{@{}p{4cm}p{10.5cm}@{}}
    \toprule
    Role & Variables \\
    \midrule
    Independent variables & Audit tenure (TENURE), audit fees (FEE), audit committee independence (ACIND) \\
    Dependent variable & Earnings management (EM), measured as absolute discretionary accruals; signed discretionary accruals (DA) in robustness checks \\
    Control variables & Firm size (SIZE), leverage (LEV), return on assets (ROA) \\
    \bottomrule
  \end{tabular}
\end{table}

\begin{table}[ht]
  \centering
  \caption{Summary of hypotheses}
  \label{tab:hypotheses}
  \small
  \begin{tabular}{@{}llp{3.5cm}@{}}
    \toprule
    \# & Hypothesis & Expected sign \\
    \midrule
    H1 & TENURE affects EM & Non-directional \\
    H2 & FEE affects EM & Non-directional \\
    H3 & ACIND affects EM & Negative \\
    \bottomrule
  \end{tabular}
\end{table}

\section{Chapter Summary}

The literature gives two competing arguments for audit tenure and for audit fees, and a clearer monitoring argument for audit committee independence. Pakistani studies tend to support lower earnings management with longer tenure and more independent committees, but they mostly pool industries and cover periods before the 2019 Regulations. Chapter~3 describes how the three hypotheses are tested in one sector under one governance regime.
